%! AP Statistics — Unit 5, Lesson 5.1 — Group Activity (Answer Key)
\documentclass[10pt]{article}
\usepackage{apstats-article}
\usepackage{apstats-key}   % pulls in -boxes; do NOT also load -boxes
\newcommand{\TallMath}[1]{$\displaystyle #1\rule[-1.4em]{0pt}{3.2em}$}

\begin{document}

\pageheader{Unit 5, Lesson 5.1}{Group Activity --- Answer Key}
\namepartnerperiod

\vspace{0.15in}

\begin{tcolorbox}[colback=redbg, colframe=black!40, boxrule=0.6pt, arc=2mm,
    left=3mm, right=3mm, top=2mm, bottom=2mm, title=\textbf{Tier R --- Remediate}]
\textbf{Directions:} Read each scenario and answer the questions. Use the sentence
frames if helpful.

\textbf{Scenario A.} A food scientist tests 30 apples from a large orchard and finds
a mean sugar content of 11.2 g. A colleague tests a different 30 apples and gets 10.8 g.

\begin{enumerate}
  \item Is either mean the true mean sugar content of all apples in the orchard?
        \begin{itemize}
          \item[$\square$] Yes \quad \ans{$\checkmark$} No
        \end{itemize}
  \item The values 11.2 and 10.8 are called \ans{statistics} because they come from a
        sample.
  \item The true mean sugar content of all apples in the orchard is called a
        \ans{parameter}.
  \item Complete: ``The two sample means differ because \ldots'' \ansline{of sampling variability.}
\end{enumerate}

\vspace{0.06in}
\textit{Sentence frame:} ``The sample means differ because of \ans{sampling variability}, which is the
natural variation in statistics from sample to sample.''
\end{tcolorbox}

\vspace{0.1in}

\begin{tcolorbox}[colback=redbg, colframe=black!40, boxrule=0.6pt, arc=2mm,
    left=3mm, right=3mm, top=2mm, bottom=2mm, title=\textbf{Tier A --- Approaching Proficiency}]
\textbf{Scenario B.} A city planner surveys 50 households in each of two
neighborhoods about daily commute time. Neighborhood 1 mean: 28.4 min.
Neighborhood 2 mean: 31.1 min.

\begin{enumerate}
  \item Identify the statistic and the parameter in each neighborhood's study.
        \ansline{Stats: 28.4 \& 31.1 min.}\\
        \ansline{Parameters: true mean commute time of all households in each.}
  \item Explain whether the 2.7-minute difference is likely due to sampling variability
        or a true population difference. What additional information would help you
        decide? \ansline{Could be chance alone.}\\
        \ansline{Samples of 50 vary, so a 2.7-min gap may be sampling variability.}\\
        \ansline{Need the SD of commute times (or a significance test) to decide.}
  \item Sketch a rough dot-plot showing what you would expect if you took 20 samples of
        size 50 from Neighborhood 1 and plotted each sample mean.

        \par\vspace{-\parskip}\nointerlineskip\noindent\begin{minipage}[t][1.0in][t]{\linewidth}\vspace*{2pt}\ans{Dot plot of 20 sample means: roughly symmetric, centered near 28.4 min,}\\
        \ans{with much less spread than individual commute times.}\end{minipage}
  \item What is the name of the distribution you just sketched? \ans{sampling distribution of $\bar{x}$}
\end{enumerate}
\end{tcolorbox}

\vspace{0.1in}

\begin{tcolorbox}[colback=redbg, colframe=black!40, boxrule=0.6pt, arc=2mm,
    left=3mm, right=3mm, top=2mm, bottom=2mm, title=\textbf{Tier E --- Extension}]
\textbf{Scenario C.} A public health researcher draws repeated random samples of size
$n = 25$ from a population whose mean daily sodium intake is $\mu = 3{,}400$ mg with
standard deviation $\sigma = 800$ mg. She records the sample mean $\bar{x}$ for each
sample.

\begin{enumerate}
  \item What is the parameter of interest? What notation describes it? \ansline{Population mean intake, $\mu$.}\\
        \ansline{$\mu = 3{,}400$ mg, the mean daily sodium intake of all people.}
  \item If she collected 10,000 sample means, predict the shape, center, and
        spread of the resulting distribution. Justify each prediction with reasoning
        (not formulas --- those come in Lesson 5.7). \ansline{Shape: symmetric.}\\
        \ansline{Center: near 3,400 mg --- sample means average out to $\mu$.}\\
        \ansline{Spread: much less than 800 mg --- averaging 25 values cancels extremes.}\\
        \ansline{So the 10,000 means pile up tightly around 3,400 mg.}
  \item How would the distribution change if she increased her sample size to $n = 100$?
        Predict the direction of change in the spread. \ansline{The spread would decrease.}\\
        \ansline{With 100 values per sample, extremes cancel out even more,}\\
        \ansline{so sample means cluster more tightly around 3,400 mg (same center).}
  \item Write one statistical question that this sampling scenario suggests about
        sodium intake in the population. \ansline{}\\
        \ansline{``Is mean daily sodium intake in our region different from 3,400 mg?''}
\end{enumerate}
\end{tcolorbox}

\end{document}
